\chapter{Software Architecture}

\section{Architectural Overview}

VirtualReactor follows a modular architecture in which chemical kinetics,
reactor physics, thermophysical properties, numerical integration, and data
management are represented by separate software components.

The architecture reflects the mathematical formulation introduced in the
previous chapter,

\begin{equation}
    \frac{\mathrm{d}\mathbf{y}}{\mathrm{d}\xi}
    =
    \mathcal{T}_{\mathrm{reactor}}
    \left(
        \mathbf{s}_{\mathrm{chemical}}
    \right)
    +
    \mathbf{s}_{\mathrm{reactor}}.
\end{equation}

Each term of this equation can be associated with a specific component of
the software architecture. This separation allows reaction mechanisms,
reactor models, and physical property models to be developed and exchanged
independently.


\section{Core Components and Initialization}

\subsection{Mechanism}

The \texttt{Mechanism} class defines the chemical reaction network and its
kinetic and thermodynamic parameters. The following reaction network is used
as an example throughout this section.

\subsubsection{Reaction Network}

Consider a system consisting of four chemical species
$A$, $B$, $C$, and $D$ connected by three reversible reactions:

\begin{align}
    A &\rightleftharpoons B
    && \text{Reaction 1}, \\
    B &\rightleftharpoons C
    && \text{Reaction 2}, \\
    2A &\rightleftharpoons D
    && \text{Reaction 3}.
\end{align}

The corresponding stoichiometric matrix is

\begin{equation}
    \boldsymbol{\nu}
    =
    \begin{bmatrix}
        -1 &  0 & -2 \\
         1 & -1 &  0 \\
         0 &  1 &  0 \\
         0 &  0 &  1
    \end{bmatrix},
\end{equation}

where each row represents a chemical species and each column represents one
reaction. Negative coefficients correspond to reactants and positive
coefficients to products.


\subsubsection{Kinetic and Thermodynamic Description}

The rate of a chemical reaction is determined by the free-energy barrier
between the reactants and the transition state. Within transition state
theory, this barrier is expressed as

\begin{equation}
    \Delta G^\ddagger(T)
    =
    \Delta H^\ddagger
    -
    T\Delta S^\ddagger,
\end{equation}

where $\Delta H^\ddagger$ and $\Delta S^\ddagger$ are the activation
enthalpy and activation entropy, respectively.

The corresponding rate constant is given by the Eyring equation,

\begin{equation}
    k(T)
    =
    \frac{k_{\mathrm{B}}T}{h}
    \exp\left(
        \frac{\Delta S^\ddagger}{R}
    \right)
    \exp\left(
        -\frac{\Delta H^\ddagger}{RT}
    \right),
\end{equation}

with the Boltzmann constant $k_{\mathrm{B}}$, the Planck constant $h$,
and the universal gas constant $R$.

The thermodynamic driving force of a reversible reaction is described
separately by the reaction Gibbs energy,

\begin{equation}
    \Delta G_{\mathrm{r}}(T)
    =
    \Delta H_{\mathrm{r}}
    -
    T\Delta S_{\mathrm{r}},
\end{equation}

which determines the equilibrium constant according to

\begin{equation}
    K(T)
    =
    \exp\left(
        -\frac{\Delta G_{\mathrm{r}}(T)}{RT}
    \right).
\end{equation}

Thus, the activation parameters determine how fast a reaction approaches
equilibrium, whereas the reaction thermodynamics determine the position of
that equilibrium.


\subsubsection{Initialization}

The reaction network can be initialized in \textsc{VirtualReactor} using the
\texttt{Mechanism} class:

\begin{verbatim}

import virtualreactor as vr


mechanism = vr.Mechanism(
    species=['A', 'B', 'C', 'D'],

    stoichiometry=np.array([
        # r1, r2, r3
        [-1,  0, -2],  # Species a
        [ 1, -1,  0],  # Species b
        [ 0,  1,  0],  # Species c
        [ 0,  0,  1]   # Species d
    ]),

    pre_exponential_factors=np.array([
        1.0e7,  # A1
        5.0e7,  # A2
        1.0e6   # A3
    ]),

    activation_energies=np.array([
        50_000.0,  # E1
        60_000.0,  # E2
        45_000.0   # E3
    ]) * ureg.joule / ureg.mol,

    reaction_enthalpies=np.array([
        -40_000.0,  # dH1: A -> B
        -30_000.0,  # dH2: B -> C
        -50_000.0   # dH3: 2 A -> D
    ]) * ureg.joule / ureg.mol,

    reaction_entropies=np.array([
        -80.0,   # dS1: A -> B
        -100.0,  # dS2: B -> C
        -150.0   # dS3: 2 A -> D
    ]) * ureg.joule / (ureg.mol * ureg.kelvin),

    reversible=np.array([
        True,
        True,
        True
    ])
)
\end{verbatim}

The order of the species and reactions must remain consistent across all
parameters. The rows of the stoichiometric matrix correspond to the species
defined in \texttt{species}, while its columns correspond to the individual
reactions. Consequently, the kinetic and thermodynamic parameter arrays
contain one entry for each reaction.

In this example, all three reactions are reversible. The arrays containing
the pre-exponential factors, activation energies, reaction enthalpies, and
reaction entropies therefore describe reactions 1--3 in the same order as
the columns of the stoichiometric matrix.



\subsection{Reactor Models}

The reactor model defines how the chemical source terms generated by the
reaction mechanism are translated into changes of the reactor state.
Different reactor models can therefore be combined with the same chemical
mechanism.

In the following example, three reactor configurations are defined:

\begin{enumerate}
    \item an ideal batch reactor,
    \item an adiabatic plug-flow reactor with constant volumetric flow, and
    \item a plug-flow reactor with external heat transfer.
\end{enumerate}

\begin{verbatim}
reactor_list = [

    vr.BatchReactor(),

    vr.PlugFlowReactor(
        flow_model=vr.ConstantVolumetricFlow(
            1.5 * ureg.liter / ureg.minute
        )
    ),

    vr.PlugFlowReactor(
        flow_model=vr.ConstantVolumetricFlow(
            1.5 * ureg.liter / ureg.minute
        ),

        jacket_temperature=300 * ureg.kelvin,

        overall_heat_transfer_coefficient=(
            300
            * ureg.watt
            / (
                ureg.meter**2
                * ureg.kelvin
            )
        ),

        heat_transfer_area_per_volume=(
            200 / ureg.meter
        )
    )
]
\end{verbatim}

The \texttt{BatchReactor} represents a closed reactor system in which the
state evolves with time. No flow model is required because no material enters
or leaves the reactor.

For a plug-flow reactor, the species state is represented by molar flow rates.
The local concentrations therefore depend on the volumetric flow rate. In this
example, the flow behavior is described by

\begin{equation}
    \dot{V}
    =
    1.5~\mathrm{L\,min^{-1}},
\end{equation}

using the \texttt{ConstantVolumetricFlow} model.

The second plug-flow reactor is adiabatic with respect to external heat
transfer, whereas the third configuration introduces heat exchange with an
external jacket. The heat-transfer contribution is characterized by the jacket
temperature

\begin{equation}
    T_{\mathrm{jacket}} = 300~\mathrm{K},
\end{equation}

the overall heat-transfer coefficient

\begin{equation}
    U = 300~\mathrm{W\,m^{-2}\,K^{-1}},
\end{equation}

and the heat-transfer area per reactor volume

\begin{equation}
    a_V = 200~\mathrm{m^{-1}}.
\end{equation}

Together, $U$ and $a_V$ determine the strength of heat exchange between the
reacting fluid and the external jacket.


\subsection{Physical Property Model}

Thermophysical properties required by the reactor energy balance are provided
independently through a \texttt{PropertyModel}. This allows the same reactor
configuration to be evaluated using different assumptions for the physical
properties of the reacting mixture.

For the present example, constant density and heat capacity are assumed:

\begin{verbatim}
constant_density_heat_capacity = vr.PropertyModel.constant(
    density=1000.0 * ureg.kg / ureg.m**3,

    heat_capacity=(
        4180.0
        * ureg.joule
        / (
            ureg.kg
            * ureg.kelvin
        )
    )
)
\end{verbatim}

The model therefore assumes

\begin{equation}
    \rho = 1000~\mathrm{kg\,m^{-3}}
\end{equation}

and

\begin{equation}
    c_p = 4180~\mathrm{J\,kg^{-1}\,K^{-1}}.
\end{equation}

These values approximately correspond to a water-like liquid and are used here
as a simple constant-property model.

Separating the physical property model from the reactor model makes it possible
to exchange property descriptions independently of the reaction mechanism or
reactor configuration. This is particularly useful for systematic parameter
studies and the generation of structured simulation datasets.


\subsection{Simulation and Initial Conditions}

Once the reaction mechanism, reactor models, and physical properties have been
defined, they are combined in a \texttt{Simulation} object. The initial state
and integration interval depend on the selected reactor model.

For the three reactor configurations introduced above, the integration limits
are defined as

\begin{verbatim}
xi_list = [300, 0.05, 0.05]
\end{verbatim}

The corresponding initial states are

\begin{verbatim}
initial_states = [
    np.array([1, 0, 0, 0, 300, 0]),      # Batch
    np.array([2.5e-5, 0, 0, 0, 300, 0]), # PFR 1
    np.array([2.5e-5, 0, 0, 0, 300, 0]), # PFR 2
]
\end{verbatim}

For the batch reactor, the state vector is

\begin{equation}
    \mathbf{y}_{0,\mathrm{Batch}}
    =
    \begin{pmatrix}
        \mathbf{c}_0 \\
        T_0 \\
        p_0
    \end{pmatrix},
\end{equation}

where the species entries represent initial concentrations. In the present
example,

\begin{equation}
    \mathbf{c}_0
    =
    \begin{pmatrix}
        1 & 0 & 0 & 0
    \end{pmatrix}^{\mathrm{T}}
    \mathrm{mol\,m^{-3}},
\end{equation}

with an initial temperature of $T_0 = 300~\mathrm{K}$.

For the plug-flow reactors, the state vector instead contains molar flow rates,

\begin{equation}
    \mathbf{y}_{0,\mathrm{PFR}}
    =
    \begin{pmatrix}
        \mathbf{F}_0 \\
        T_0 \\
        p_0
    \end{pmatrix},
\end{equation}

with

\begin{equation}
    \mathbf{F}_0
    =
    \begin{pmatrix}
        2.5\times10^{-5} & 0 & 0 & 0
    \end{pmatrix}^{\mathrm{T}}
    \mathrm{mol\,s^{-1}}.
\end{equation}

The independent coordinate $\xi$ is interpreted according to the reactor
model. For the batch reactor,

\begin{equation}
    \xi = t,
\end{equation}

and the simulation is performed over

\begin{equation}
    0 \leq t \leq 300~\mathrm{s}.
\end{equation}

For the plug-flow reactors,

\begin{equation}
    \xi = V,
\end{equation}

and the reactor state is integrated over

\begin{equation}
    0 \leq V \leq 0.05~\mathrm{m^3}.
\end{equation}


\subsubsection{Running the Simulations}

The simulations are performed iteratively using the same reaction mechanism
and physical property model while exchanging the reactor configuration:

\begin{verbatim}
results = []

for i in range(3):

    vdv_sim = vr.Simulation(
        mechanism=van_de_vusse,
        reactor=reactor_list[i],
        properties=constant_density_heat_capacity
    )

    vdv_sim_result = vdv_sim.solve(
        initial_state=initial_states[i],
        xi_span=(0, xi_list[i])
    )

    results.append(vdv_sim_result)
\end{verbatim}

Each \texttt{Simulation} combines three independent model components:

\begin{equation}
    \boxed{
    \texttt{Simulation}
    =
    \texttt{Mechanism}
    +
    \texttt{Reactor}
    +
    \texttt{PropertyModel}
    }
\end{equation}

The \texttt{solve()} method integrates the reactor state from the specified
initial condition over the interval \texttt{xi\_span}. The resulting
\texttt{SimulationResult} is appended to \texttt{results}, allowing the
different reactor configurations to be analyzed and compared within the same
workflow.

This modular initialization makes it possible to systematically vary reactor
models, operating conditions, or physical property models while keeping the
chemical mechanism unchanged.


\subsubsection{Simulation Results}

The \texttt{solve()} method performs the numerical integration and returns
a \texttt{SimulationResult} object,

\begin{equation}
    \texttt{Simulation}
    \xrightarrow{\;\texttt{solve()}\;}
    \texttt{SimulationResult}.
\end{equation}

A \texttt{SimulationResult} represents the complete output of a single
simulation. It stores the values of the independent coordinate $\xi$ together
with the corresponding reactor states $\mathbf{y}(\xi)$.

The result of a simulation can therefore be obtained as

\begin{verbatim}
result = simulation.solve(
    initial_state=initial_state,
    xi_span=(0, xi_end)
)
\end{verbatim}

Conceptually, the numerical solution consists of

\begin{equation}
    \left\{
        \xi_i,\,
        \mathbf{y}(\xi_i)
    \right\}_{i=1}^{N},
\end{equation}

where $N$ is the number of evaluated points along the simulation coordinate.

For a batch reactor, the stored solution corresponds to

\begin{equation}
    \left\{
        t_i,\,
        \mathbf{c}(t_i),\,
        T(t_i),\,
        p(t_i)
    \right\},
\end{equation}

whereas for a plug-flow reactor it corresponds to

\begin{equation}
    \left\{
        V_i,\,
        \mathbf{F}(V_i),\,
        T(V_i),\,
        p(V_i)
    \right\}.
\end{equation}

Separating the numerical result from the \texttt{Simulation} object allows
simulation data to be analyzed, plotted, compared, and stored independently
of the numerical integration itself.


\subsubsection{Serialization}

Simulation results can be serialized for persistent storage and subsequent
analysis. \textsc{VirtualReactor} uses the HDF5 format to store simulation
data together with the information required to describe the corresponding
simulation.

A result can be written to an HDF5 file using

\begin{verbatim}
result.save_hdf5(
    "data/raw/pfr_simulation.h5"
)
\end{verbatim}

Serialization converts the in-memory representation of a
\texttt{SimulationResult} into a persistent file representation,

\begin{equation}
    \texttt{SimulationResult}
    \xrightarrow{\;\texttt{save\_hdf5()}\;}
    \texttt{HDF5 file}.
\end{equation}

The HDF5 format is particularly suitable for numerical simulation data because
it can efficiently store multidimensional numerical arrays together with
structured metadata. This allows simulation results to be retained in a
reproducible and machine-readable form and provides a basis for systematic
parameter studies and the generation of larger simulation datasets.


\subsubsection{Serialization}

A \texttt{SimulationResult} exists initially as a Python object in memory.
For persistent storage, the result can be serialized into an HDF5 file using

\begin{verbatim}
result.save_hdf5(
    "simulation.h5"
)
\end{verbatim}

The resulting file preserves both the numerical solution and the information
required to interpret the stored state. Its hierarchical structure can be
represented as

\begin{verbatim}
simulation.h5
|
+-- coordinates
+-- states
+-- species
|
+-- attributes
    +-- file_format
    +-- schema_version
    +-- coordinate_name
    +-- species_state_name
    +-- success
    +-- message
\end{verbatim}

The \texttt{coordinates} dataset contains the values of the independent
simulation coordinate $\xi$. The \texttt{states} dataset contains the
corresponding reactor states $\mathbf{y}(\xi)$, while \texttt{species}
stores the associated species names.

Additional information is stored as HDF5 attributes. In particular,
\texttt{coordinate\_name} identifies the independent coordinate, such as
time or reactor volume, and \texttt{species\_state\_name} identifies whether
the species states represent concentrations or molar flow rates.

A serialized result can subsequently be reconstructed using

\begin{verbatim}
result = vr.SimulationResult.load_hdf5(
    "simulation.h5"
)
\end{verbatim}
